\section{Первая игра: от заготовки до стола}\label{tut:first-session}

\paragraph{Цель}
Провести короткую игру по сценарию из раздела \ref{tut:first-scenario} --- с замком, дворецким и запретной дверью. Разбор идёт по шагам: что нажать и что при этом увидят игроки.

\paragraph{Что нужно перед началом}
Готовый сценарий с локацией ``Замок'', NPC ``Дворецкий'' и story beat ``Встреча в замке''. Один-два игрока с аккаунтами на сайте.

\subsection{Шаг 1. Запустить мир}\label{tut:first-session:launch}

Откройте сценарий и переведите его в запущенное состояние. Появится запущенный мир --- он попадёт в список \texttt{/launched-scenarios}.

С этого момента сценарий и мир живут отдельно: урон, потраченные ресурсы и открытые области карты меняют мир, а заготовка остаётся нетронутой. Захотите провести ту же историю для другой компании --- запустите сценарий ещё раз.

\subsection{Шаг 2. Собрать игроков}\label{tut:first-session:lobby}

Создайте лобби и отправьте игрокам ссылку. Быстрый вариант --- команда в чате партии (см.\ раздел \ref{sec:telegram:lobby}): бот разошлёт каждому персональную ссылку сразу в лобби.

Дождитесь, пока все подключатся, и \textbf{назначьте каждому персонажа}. Не пропускайте этот шаг: игрок без персонажа попадёт в игру и увидит пустой экран.

\subsection{Шаг 3. Открыть первую сцену}\label{tut:first-session:first-scene}

Запустите игру --- всех перебросит в сессию.

Выберите локацию ``Замок'' и откройте story beat ``Встреча в замке''. Заготовленный текст сцены уйдёт игрокам: они прочитают про холл, дворецкого и запрет приближаться к двери.

Представьте дворецкого: покажите его карточку всем участникам. NPC выведется крупно на экранах игроков --- нагляднее, чем описывать словами.

\subsection{Шаг 4. Первая проверка}\label{tut:first-session:first-roll}

Допустим, игрок решает уговорить дворецкого пропустить его к двери. В сценарии для этого заготовлено препятствие ``Убедить дворецкого'' со сложностью 12.

Игрок объявляет соответствующий ход и указывает цель --- дворецкого. Дальше:

\begin{enumerate}
\item другие игроки могут помочь --- их вклад добавится к броску;
\item система бросает кости, результат появляется в ленте у всех;
\item по результату определяется исход: успех, частичный успех или провал;
\item вы выбираете конкретное последствие и подтверждаете изменения.
\end{enumerate}

Частичный успех --- самый интересный случай: дворецкий отвлекается, но зовёт охрану. Заготовленная формулировка последствий очень выручает именно здесь, когда придумывать некогда.

\subsection{Шаг 5. Развитие сцены}\label{tut:first-session:develop}

Пока группа действует, вам пригодятся:

\begin{itemize}
\item \textbf{счётчик} --- заведите ``Подозрительность дворецкого'' и увеличивайте на каждой неудаче. Игроки видят рост шкалы, напряжение нарастает само;
\item \textbf{личное сообщение} --- отправьте одному игроку ``ты замечаешь ключ на поясе дворецкого''. Остальные этого не увидят;
\item \textbf{области карты} --- если у замка есть карта, открывайте комнаты по мере продвижения группы;
\item \textbf{музыка} --- поставьте трек, он заиграет у всех одновременно.
\end{itemize}

\subsection{Шаг 6. Завершить подход}\label{tut:first-session:finish}

Когда встреча закончена, завершите подход. Проверьте перед этим, что все изменения внесены: \textbf{вернуть завершённый подход нельзя}.

Если мир входит в кампанию, состояние персонажей, инвентарь и счётчики перенесутся в следующую встречу. Если нет --- следующий подход начнётся из текущего состояния мира, но без переноса состояния персонажей.

\paragraph{Что дальше}
Мир остался жив и ждёт следующей игры: дверь всё ещё заперта, дворецкий всё помнит, счётчик подозрительности сохранил значение. Следующую встречу начинайте с шага \ref{tut:first-session:lobby}.
